A contribuição de \citep{zagirov2023modeling}  é significativa para o campo da robótica pois propõe um modelo virtual do robô humanoide AR601M no simulador Gazebo. Estudos relacionados ao AR601M exploraram aspectos de locomoção bípede, algoritmos de estabilidade e otimização de padrões de marcha, como destacado por \citep{zagirov2023modeling}. Esses trabalhos abordaram métodos cinemáticos e dinâmicos para aprimorar a mobilidade do robô, com ênfase no controle de equilíbrio durante a caminhada. Outros estudos também avançaram no controle das mãos do robô, implementando algoritmos de preensão e manipulação de objetos em ambientes simulados.

Embora relevantes, essas pesquisas individuais eram limitadas em integrar as habilidades motoras grossas e finas em um único modelo funcional. A proposta de \citep{zagirov2023modeling} distingue-se ao consolidar esses avanços em um modelo abrangente, permitindo simular ações humanas como caminhar, movimentar os braços e manipular objetos em ambientes colaborativos. Ao contrário dos modelos padrão do Gazebo, o AR601M possui maior capacidade de interação física com o ambiente e os sensores, superando limitações anteriores. Esse enfoque representa um progresso significativo na modelagem de interações humano-robô, com implicações diretas para a segurança e a eficiência em espaços colaborativos.

Outros estudos destacam a relevância dos simuladores para experimentos em robótica espacial, permitindo validar algoritmos antes da aplicação física. \citep{koenig2004design} ressaltaram o potencial do Gazebo como motor de simulação para sistemas robóticos complexos, enquanto \citep{quigley2009ros} demonstraram como o ROS facilita a integração e o controle de robôs em simulações. Trabalhos como os de \citep{smith2016astrobee} e \citep{bualat2018astrobee}, exploraram o uso de simuladores para robôs espaciais como o Astrobee e o VIPER, reforçando a importância de ambientes realistas para testar sistemas de navegação e manipulação. Nesse contexto, \citep{santana2021development} avançaram ao desenvolver um simulador de alta fidelidade no Gazebo, integrando sensores, propulsão e algoritmos de visão computacional para testar operações de proximidade no espaço, contribuindo significativamente para o campo.

